
%\documentclass[oneside,final,14pt]{extreport}
\documentclass[12pt, a4paper, openany]{book}



\usepackage[left=2cm,right=2cm,top=2cm,bottom=2cm,bindingoffset=0cm]{geometry}
%\usepackage[koi8-r]{inputenc}
%\usepackage[russianb]{babel}
\usepackage{vmargin}

\setpapersize{A4}
\usepackage[T2A]{fontenc}
\usepackage[utf8x]{inputenc}
\usepackage[english, russian]{babel}
\setmarginsrb{2cm}{2cm}{2cm}{2cm}{0pt}{5mm}{0pt}{0mm}
\usepackage{indentfirst}
\usepackage{nicefrac} % For comparison
%\usepackage{xfrac} % Works better with other fonts
%\usepackage[unicode, pdftex]{hyperref}
\usepackage{lettrine}
\usepackage[usenames]{color}
\usepackage{colortbl}
\usepackage{mathtext}
\usepackage{epigraph}
\usepackage{amsmath, amsfonts, amssymb, mathrsfs}
%\usepackage{mathptmx}
%\usepackage{txfonts}
\usepackage{pxfonts}
\usepackage[pagestyles]{titlesec}
\usepackage{ebgaramond}
\usepackage{awesomebox}
\usepackage{enumitem}
\usepackage{makeidx}
\usepackage[letterspace=150]{microtype}
\makeindex
\usepackage{etoolbox}
\makeatletter
\newlength\epitextskip
\pretocmd{\@epitext}{\em}{}{}
\apptocmd{\@epitext}{\em}{}{}
\patchcmd{\epigraph}{\@epitext{#1}\\}{\@epitext{#1}\\[\epitextskip]}{}{}
\makeatother
%running fraction with slash - requires math mode.
\newcommand*\rfrac[2]{{}^{#1}\!/_{#2}}

\DeclareSymbolFont{Xlargesymbols}{OMX}{cmex}{m}{n}
\DeclareMathSymbol{\Xsum}{\mathop}{Xlargesymbols}{80}

\setlength\epigraphrule{0pt}
\setlength\epitextskip{2ex}
\setlength\epigraphwidth{.8\textwidth}

\usepackage{xfrac} % Works better with other fonts
\usepackage[colorlinks=true,linkcolor=black,urlcolor=black,bookmarksopen=true]{hyperref}

\usepackage{fancyhdr} % пакет для установки колонтитулов
\pagestyle{fancy} % смена стиля оформления страниц
\fancyhf{} % очистка текущих значений
\fancyhead[C]{\thepage} % установка верхнего колонтитула
\renewcommand{\headrulewidth}{0pt} % убрать разделительную линию


% Настройка вертикальных и горизонтальных отступов
\titlespacing{\chapter}{0pt}{5pt}{5pt}
\titlespacing{\section}{\parindent}{4mm}{4mm}
\titlespacing{\subsection}{\parindent}{3mm}{3mm}


%
%\renewcommand{\tabcolsep}{1cm} %% increase table column spacing

% Настройка задачи со зведочкой
\newcounter{namedlistcounter} % number the items
\newenvironment{withdot}
{\begin{list}
		{\arabic{namedlistcounter}*.} % labeling 
		{\usecounter{namedlistcounter} % set counter
			\setlength{\leftmargin}{3em}} % set spacing 
	}
	{\end{list}}


\newcommand{\anonsection}[1]{ \section*{#1} \addcontentsline{toc}{section}{\numberline {}#1}} 

\makeatletter %%%%% <---- Starting chapter without a pagebreak
\renewcommand\chapter{\par%
	\thispagestyle{plain}% \global\@topnum\z@
	\@afterindentfalse \secdef\@chapter\@schapter}
\makeatother %%%%% <---- Starting chapter without a pagebreak
\titleformat{\chapter}[display]
{\normalfont\bfseries}{}{0pt}{\Large}

\newpagestyle{mystyle}{
	\sethead[\thepage][][]{}{}{\thepage}	
}

\renewcommand{\rmdefault}{cmr}


\pagestyle{mystyle}
\sloppy
\begin{document}

	
	\begin{titlepage}
		
		\begin{center}
			\topskip0pt
			\vspace*{\fill}
			
			
			{\large\bf Г. АЛЕКСАНДРОВ\\}
			\ \\
			\ \\
			{\Huge\bf БИБЛИЯ ЛЮДОЕДОВ\\}
			\ \\
			\ \\
			\ \\
			\vspace*{\fill}
			
			\vfill
			
			
			О Г И З
			
			ГОСУДАРСТВЕННОЕ ИЗДАТЕЛЬСТВО
			
			ПОЛИТИЧЕСКОЙ ЛИТЕРАТУРЫ
			
			1\ 9\ 4\ 2
			

		\end{center}
		
	\end{titlepage}
	
	\thispagestyle{empty} % выключаем отображение номера для этой страницы
	
	\newpage
	

\setcounter{secnumdepth}{0}
	
	\section[Фашизм --- лютый враг человечества]{\center \textbf{ФАШИЗМ --- ЛЮТЫЙ ВРАГ ЧЕЛОВЕЧЕСТВА}}
	
	\subsection*{\center \textbf{Кто правит нынейшей Германией}}

Жалкие выродки истории, как метко назвал немецких фашистов Горький, завсегдатаи кабаков и участники почти всех преступлений, совершенных за последние годы в Германии, разбойничьи захватив власть и назвав себя всем на посмешище «национал-социалистами», вот уже несколько лет организуют одну кровавую бойню за другой.

Германия не переживала еще столь позорных страниц своей истории, как в годы бандитской власти шайки разнузданных авантюристов — гитлеровских головорезов.

Страна, давшая миру Гете и Лессинга, Гейне и Бетховена, Гегеля и Маркса, теперь вынуждена под страхом штыка и виселицы до поры до времени молчаливо терпеть, переносить кровавые оргии, авантюристические выходки, преступные действия своих столь же наглых, сколь и недальновидных правителей. Очутившиеся у власти гитлеровские преступники вовлекли в войну большинство населения земного шара, вероломно порвали свои собственные договоры, уже уничтожили и продолжают уничтожать сотни тысяч самой молодой и сильной части немецкого народа, вырезают, истребляют все передовое, честное, свободолюбивое в немецком народе, ввергая его, а также и народы европейских стран в одну кровавую трясину за другой.

Кто же эти зарвавшиеся авантюристы, кровожадные псы с клеймом свастики, не признающие никакого права, кроме нагло применяемой силы для восстановления рабства? Кто эти люди, для которых все продажно, нет ничего святого, которые в любую минуту готовы на самые ужасные кровавые преступления для удовлетворения своих столь же корыстных, сколь и преступных инстинктов и замыслов?

Люди, которых проклинают и будут проклинать миллионы, имена которых народы всегда будут произносить с ненавистью и чувством величайшего омерзения, — Гитлер, Геринг, Геббельс, Фрик, Риббентроп и другие — это выродки с темным и разнообразным прошлым. Но одно их объединяет — злобная, лютая ненависть к народным массам как самой Германии, так и других стран. Сами фашисты пишут, впрочем, об этом. Так, в изданной недавно книге фашиста Ганса Фрейера открыто и официально утверждается, что, «опираясь на всякий ‹брод, политический гений строит блеск своего господства».

Для правящей клики фашистской Германии нет страшнее врага, чем свой собственный народ, — она его боится ий ненавидит. Политические авантюристы и бандиты — они не могут иначе относиться к народу: по Гитлеру, «понимание масс ограничено, их разум мал», народ «сам по себе ленив»; главный мастер Гитлера по делам демагогии и оглупления масс — Геббельс призывал «снова обратиться к самым примитивным инстинктам масс».

Разрушение, уничтожение, истребление! Концлагерь, виселица, расстрел! —вот наиболее часто встречающиеся слова в речах и статьях правителей нынешней Германии, когда они обращаются к народу. Эти люди ненавидят все новое, передовое. Такие понятия, как прогресс, свобода народов, вызывают у них тошноту, ужас и злобу. Ограниченные, тупые и вместе с тем ловкие и наглые демагоги, они жаждут крови и интересуются лишь тем, что способствует возникновению новых войн, новых кровавых столкновений народов. Война для них — «нормальное» состояние человечества. «Жизнь есть война» — это крылатое выражение стремится вдолбить в солдатские головы немецкой армии ее командование, ибо мирный труд означает смерть кровавого фашистского чудовища; в смерти же, в гибели рабочих, крестьян, интеллигентов на войне фашистские правители видят средство сохранения своей разбойничьей, воровской, антинародной власти. Недаром Гитлер с циничностью наемного убийцы сказал однажды: «Война — это я».

Разве не ясно, что преступная гитлеровская шайка, пришедшая к власти через тысячи трупов немецких рабочих, крестьян, служащих, ученых, писателей, сделавшая бандитские приемы единственным орудием своей политики, зарвавшаяся в разбойничьих военных авантюрах, создавшая правительство, состоящее из провокаторов, шпионов, «пивных политиков», палачей и насильников, не может победить, но будет с беспощадностью разгромлена и уничтожена!

	\subsection*{\center \textbf{Фашизм --- злейший враг человечества}}


Фашизм — злейший враг всего передового, прогрессивного человечества.

Фашизм, гитлеризм есть открытая террористическая диктатура наиболее реакционных, наиболее шовинистических и наиболее империалистических элементов финансового капитала.

Фашизм — главный враг всех народов и государств, ибо он несет им кабалу, рабство, искоренение. Фашистское движение — лютый враг всех народов. Подонки современного общественного строя в Германии и Италии, фашисты рассматривают всех соучастников своих кровавых преступлений как «высшую расу», призванную покорять чужие народы — якобы «низшую расу» — и жить за их счет. Фашизм несет на штыках своей армии уничтожение всех достижений современной цивилизации, всех культурных завоеваний и свобод, достигнутых народами. Фашизм — враг народов; он несет им, как показывает положение в оккупированных Германией и Италией странах, уничтожение национальной самостоятельности и целостности народов к государства, ликвидацию промышленности, ограбление всех запасов продовольствия, а следовательно, голод и смерть сотням тысяч, миллионам людей.

Истинные цели фашистских правителей ясны из той политики, которая проводится ими в оккупированных странах — во Франции, Чехословакии, Польше, Норвегии, Бельгии, Дании, Голландии, Греции, Югославии, Румынии и других странах. Фашистские войска погружают в этих странах на транспорт и отправляют в Германию конфискованное у населения продовольствие и предметы первой необходимости, вывозят имеющиеся в стране сырье, машины, оборудование, сгоняют тысячи мирного населения на немецкие поля и под плетью и штыком солдат заставляют, как в древности рабов, днем и ночью нести подневольный, унизительный труд; голодных и обессилевших людей при первом даже робком протесте пристреливают, на утеху своему «фюреру» вешают на первых же деревьях или сгоняют в вонючие ямы, называемые концентрационными лагерями.

Гнев народов оккупированных и завоеванных стран — это страшная для фашистских правителей сила. История знает многочисленные примеры, когда гневное возмущение покоренных, но свободолюбивых народов могучим порывом сбрасывало иго иностранных поработителей. А причин для нарастания такого народного негодования больше чем достаточно.

В самом деле, подавляющее большинство французов, сжав кулаки, страстно желая полного разгрома современных варваров, наблюдает, как фашисты хозяйничают во Франции; голландцы, бельгийцы, датчане, норвежцы живут с мыслью о том, чтобы отомстить фашистской Германии. Многие десятки тысяч рабочих, насильно увезенных из Бельгии и Дании в Германию для рабской работы; изъятие по смехотворным, бросовым ценам, установленным немцами, всех сельскохозяйственных продуктов, что уже вызвало массовое разорение крестьянства в этих странах; назначение администраторами школ и университетов оккупированных стран фашистских головорезов — все это не могло не вызвать чувства огромного возмущения и ненависти к поработителям.

Никогда не забудет всех своих мук и страданий польский народ. На глазах всего населения фашистами были зверски убиты многие сотни тысяч поляков; десятки тысяч польских семей оккупанты выгнали из родных домов, отдав немцам все их имущество; в польских городах и поселках до сих пор устраиваются облавы на девушек, которых насилуют и сотнями отправляют в дома терпимости.

Мало чем отличается положение и в других оккупированных странах — Чехии и Югославии, Греции и Румынии.

Как же после этого завоеванным народам Европы страстно не желать уничтожения фашизма, подлой гитлеровской своры! Как же им не стремиться всеми средствами и силами содействовать поражению фашизма в нынешней войне!

Разве не ясно, что фашизм, провозглашающий покорение и закабаление других народов своей целью, рабство — своим высшим принципом, человеконенавистничество и кровожадность — своим идеалом, не может победить, но неминуемо будет уничтожен!

Разве не ясно, что шайка фашистских извергов — нынешние правители Германии, — навлекшая на немецкий народ
	
	
	\newpage
	\tableofcontents
	
	\thispagestyle{empty} % 
	
	\newpage
	
	\setcounter{secnumdepth}{0}
	
	\phantomsection
	
		\section*{Описание}
	
	{\bf Название:} А если в семье телевизор... и дети?
	
{\bf Автор:} Ксенофонтов Валентин Васильевич

{\bf Редактор:} О. Г. Свердлова
	
{\bf Издательство:} М.: <<Знание>>. 1973. - 96 с.
	
		{\bf Аннотация:} Ученые всех стран мира сходятся в одном — обилие телевизионной информации, даже полезной по своему содержанию, без правильного руководства со стороны взрослых воспитывает у детей душевную и нравственную пассивность, приучает довольствоваться малым, не говоря уже о том, что оно отвлекает от чтения, спорта и других занятий. И в то же время телевидение представляет детям весь мир наглядно, в ярких динамичных образных формах. Чего только не увидишь по телевизору! Луноход и дно океана, сказку и экзотических животных. 
		
		Книга учит родителей правильно использовать огромные возможности телевидения для воспитания детей. Она рассчитана на самые широкие круги читателей. 
		


		\thispagestyle{empty} % выключаем отображение номера для этой страницы

	
\end{document}


